\documentclass[aspectratio=169,10pt]{beamer}
\usepackage[T1]{fontenc}
\usepackage{lmodern}
\usepackage{listings}
\usepackage{xcolor}

\definecolor{Ink}{HTML}{172B3A}
\definecolor{Blue}{HTML}{235789}
\definecolor{Muted}{HTML}{526574}
\definecolor{Paper}{HTML}{F5F7F8}
\definecolor{Code}{HTML}{EAF0F4}
\setbeamercolor{normal text}{fg=Ink,bg=white}
\setbeamercolor{frametitle}{fg=Ink,bg=white}
\setbeamercolor{title}{fg=Ink}
\setbeamercolor{structure}{fg=Blue}
\setbeamercolor{block title}{fg=Blue,bg=white}
\setbeamercolor{block body}{fg=Ink,bg=Paper}
\setbeamertemplate{navigation symbols}{}
\setbeamertemplate{footline}[frame number]
\setbeamersize{text margin left=0.7cm,text margin right=0.7cm}
\setlength{\parskip}{0.35em}
\lstset{basicstyle=\ttfamily\small,columns=fullflexible,keepspaces=true,
  showstringspaces=false,breaklines=true,backgroundcolor=\color{Code},
  frame=single,rulecolor=\color{Code},xleftmargin=0.4em,xrightmargin=0.4em,
  aboveskip=0.6em,belowskip=0.6em}

\title{Continuation-Passing Style}
\subtitle{Making the rest of the computation explicit}
\author{CS413}
\date{Fall 2026}

\begin{document}
\begin{frame}[plain]\titlepage\end{frame}

% 2 minutes
\begin{frame}{A computation has a result and a next step}
  Consider evaluating \texttt{(2 + 3) * 4}. After computing \texttt{2 + 3},
  the program still has work to do: multiply that result by 4.
  \begin{center}
    \Large \texttt{2 + 3} \quad then \quad \texttt{multiply by 4}
  \end{center}
  \begin{block}{Continuation}
    A continuation is a function that says what to do with a value when that
    value becomes available.
  \end{block}
  Here the continuation for the sum is \texttt{fun x -> x * 4}. The final
  answer is obtained by giving it the value \texttt{5}.
\end{frame}

% 2 minutes
\begin{frame}[fragile]{Direct style hides the next step}
  \begin{lstlisting}
fun squarePlusOne(n) = n * n + 1
  \end{lstlisting}
  A direct-style function returns its answer to its caller. The caller's
  remaining work sits outside the function.
  \begin{center}
    \texttt{squarePlusOne(3)} \quad returns \quad \texttt{10}
  \end{center}
  \vspace{0.5em}
  CPS changes the interface: instead of returning a result, a function
  receives a continuation and passes its result to that continuation.
\end{frame}

% 3 minutes
\begin{frame}[fragile]{A function in CPS receives its continuation}
  \begin{columns}[T]
    \column{.48\textwidth}
    \textbf{Direct style}
    \begin{lstlisting}
squarePlusOne(n) = n*n + 1
    \end{lstlisting}
    Result type: \texttt{int}
    \column{.48\textwidth}
    \textbf{Continuation-passing style}
    \begin{lstlisting}
squarePlusOne_cps(n, k) =
  k(n*n + 1)
    \end{lstlisting}
    Result type: \texttt{(int -> R) -> R}
  \end{columns}
  \begin{block}{The identity continuation}
    \texttt{fun x -> x} asks for the final value unchanged. Passing it in
    recovers the ordinary result: \texttt{squarePlusOne\_cps(3, id) = 10}.
  \end{block}
  The type variable \texttt{R} describes the eventual answer produced by the
  continuation; it need not be an integer.
\end{frame}

% 3 minutes
\begin{frame}[fragile]{Translate an expression by naming its next step}
  To translate \texttt{e} in a context that expects its value, write
  \texttt{CPS(e, k)}.
  \begin{lstlisting}
CPS(42, k) = k(42)
CPS(x,  k) = k(x)
  \end{lstlisting}
  For an addition, evaluate the left operand, then the right, then continue:
  \begin{lstlisting}
CPS(e1 + e2, k) =
  CPS(e1, fun v1 ->
    CPS(e2, fun v2 ->
      k(v1 + v2)))
  \end{lstlisting}
  The nested functions remember the unfinished work. This translation makes
  left-to-right evaluation order explicit.
\end{frame}

% 3 minutes
\begin{frame}[fragile]{Example: CPS for a nested arithmetic expression}
  Translate \texttt{(a + b) * c}:
  \begin{lstlisting}
CPS(a + b, fun sum ->
  CPS(c, fun factor ->
    k(sum * factor)))
  \end{lstlisting}
  \begin{enumerate}
    \item Evaluate \texttt{a + b}; call the continuation with \texttt{sum}.
    \item Evaluate \texttt{c}; call the next continuation with \texttt{factor}.
    \item Multiply and hand the answer to \texttt{k}.
  \end{enumerate}
  With \texttt{a=2}, \texttt{b=3}, \texttt{c=4}, and identity \texttt{k},
  the result is \texttt{20}.
\end{frame}

% 4 minutes
\begin{frame}[fragile]{Recursive calls must also pass the rest of the work}
  In factorial, after the recursive call returns, we still multiply by \texttt{n}.
  CPS stores that multiplication in a continuation:
  \begin{lstlisting}
fact_cps(n, k) =
  if n <= 1 then k(1)
  else fact_cps(n - 1,
    fun r -> k(n * r))
  \end{lstlisting}
  \begin{itemize}
    \item Base case: the answer is 1, so pass 1 to \texttt{k}.
    \item Recursive case: when \texttt{fact(n-1)} produces \texttt{r}, multiply by \texttt{n}, then continue.
    \item Start with \texttt{fact\_cps(4, fun x -> x)}; the result is 24.
  \end{itemize}
\end{frame}

% 3 minutes
\begin{frame}{The continuation records pending multiplications}
  For \texttt{fact\_cps(4, id)}, the recursive descent builds this chain:
  \begin{center}
    \texttt{id} $\leftarrow$ \texttt{fun r -> id(4*r)} $\leftarrow$
    \texttt{fun r -> k4(3*r)} $\leftarrow$ $\cdots$
  \end{center}
  \vspace{0.6em}
  The full continuation eventually represents
  \[
    \texttt{fun r -> id(4 * (3 * (2 * (1 * r))))}
  \]
  \begin{block}{What changed?}
    Direct style leaves pending work in the call stack. CPS represents that
    work as functions passed explicitly from call to call.
  \end{block}
  \small (The displayed chain is schematic: each continuation closes over
  the values needed for its multiplication.)
\end{frame}

% 3 minutes
\begin{frame}[fragile]{Conditionals and function calls have CPS rules too}
  \begin{lstlisting}
CPS(if test then yes else no, k) =
  CPS(test, fun b ->
    if b then CPS(yes, k)
         else CPS(no, k))
  \end{lstlisting}
  Only the chosen branch runs, and both branches receive the same continuation.
  \vspace{0.5em}
  A call \texttt{f(x)} first computes the function and argument, then applies
  the function to the argument and continuation:
  \begin{lstlisting}
CPS(f(x), k) = CPS(f, fun vf ->
  CPS(x, fun vx -> vf(vx, k)))
  \end{lstlisting}
  The exact rule depends on the language's evaluation order and calling convention.
\end{frame}

% 2 minutes
\begin{frame}{CPS changes function types}
  A direct function returning \texttt{A} and accepting \texttt{B} has type
  \[
    \texttt{B -> A}
  \]
  Its CPS version accepts a continuation from \texttt{A} to some answer type \texttt{R}:
  \[
    \texttt{B -> (A -> R) -> R}
  \]
  \begin{itemize}
    \item \texttt{A}: the value the direct function would have returned.
    \item \texttt{R}: the answer type chosen by the surrounding computation.
  \end{itemize}
  The continuation argument adds another function layer; polymorphic CPS may
  quantify over \texttt{R}.
\end{frame}

% 2 minutes
\begin{frame}{CPS and tail calls}
  In \texttt{fact\_cps}, the recursive call is the final action in its branch:
  \begin{center}
    \texttt{fact\_cps(n - 1, continuation)}
  \end{center}
  It is a tail call: the caller has no instructions left after it returns.
  \begin{itemize}
    \item Tail-call execution can reuse the current stack frame.
    \item CPS makes pending work explicit, but a continuation may still be a closure that uses memory.
    \item CPS by itself does not guarantee constant space; representation and runtime behavior matter.
  \end{itemize}
\end{frame}

% 2 minutes
\begin{frame}{CPS can represent more than ordinary returns}
  A function can choose which continuation to invoke:
  \begin{itemize}
    \item \textbf{Success and failure:} pass separate continuations for each outcome.
    \item \textbf{Exceptions:} send errors to an exceptional continuation.
    \item \textbf{Early exit:} invoke a continuation that skips the remaining work.
    \item \textbf{Control operators:} capture and later invoke a saved continuation.
  \end{itemize}
  These ideas all make control flow visible as data: functions that can be
  passed, stored, or selected.
\end{frame}

% 1 minute
\begin{frame}{Summary}
  \begin{itemize}
    \item A continuation describes the rest of a computation.
    \item CPS passes that continuation explicitly instead of returning a value.
    \item Translation nests continuations around unfinished operations.
    \item CPS exposes control flow and makes calls tail calls, while closures may still consume space.
  \end{itemize}
  \vspace{1em}
  \begin{block}{Quick check}
    In \texttt{fact\_cps}, why does the recursive call receive
    \texttt{fun r -> k(n*r)} instead of \texttt{k} directly?
  \end{block}
\end{frame}

% Optional example
\begin{frame}[fragile]{Optional example: Fibonacci in direct style}
  The Fibonacci function makes two recursive calls. The second call and the
  addition are work that remains after the first call returns:
  \begin{lstlisting}
fib(n) =
  if n <= 1 then n
  else fib(n - 1) + fib(n - 2)
  \end{lstlisting}
  For example, \texttt{fib(4)} computes \texttt{fib(3) + fib(2)}, or
  \texttt{2 + 1 = 3}. This definition returns its result normally; the caller
  does not see the pending second call or addition.
\end{frame}

\begin{frame}[fragile]{Optional example: Fibonacci in CPS}
  After \texttt{fib(n-1)} produces \texttt{left}, the function still needs
  to compute \texttt{fib(n-2)} and add the two results. The continuation stores
  that work:
  \begin{lstlisting}
fib_cps(n, k) =
  if n <= 1 then k(n)
  else fib_cps(n - 1, fun left ->
    fib_cps(n - 2, fun right ->
      k(left + right)))
  \end{lstlisting}
  Both recursive calls pass their results onward through continuations. With
  the identity continuation, \texttt{fib\_cps(4, fun x -> x)} evaluates to 3.
\end{frame}
\begin{frame}[fragile]{Optional example: Fibonacci as lazy call events}
  Record both the entry to a call and its return. The continuation produces
  the remaining event stream after a recursive result is available:
  \begin{lstlisting}
event = Enter(n) | Return(value)
emit(event, rest) = lazy_cons(event, delay(rest))
fib_events(n, k) =
  emit(Enter(n),
    if n <= 1 then emit(Return(n), k(n))
    else fib_events(n - 1, fun left ->
      fib_events(n - 2, fun right ->
        emit(Return(left + right), k(left + right)))))
  \end{lstlisting}
  Start with \texttt{fib\_events(4, fun result -> lazy\_empty)}. The outer
  call emits \texttt{Enter(4)}; its recursive calls emit their own entry and
  return events before the outer call emits \texttt{Return(3)}. The tail is
  delayed, so events are produced as the consumer advances through the stream.
\end{frame}

\begin{frame}{The event consumer drives the computation}
  Each invocation emits \texttt{Enter(n)} before its work and
  \texttt{Return(value)} after its work. The continuation keeps the pending
  second recursive call and addition, so returns appear in evaluation order.
  \begin{block}{Event stream}
    These are the only two event forms. The final event is the outer call's
    return, whose value is the Fibonacci result.
  \end{block}
  A real implementation needs a lazy-list type whose tail is suspended, such
  as a thunk or stream node.
\end{frame}
\end{document}
